\documentclass{article}
\usepackage[T1]{fontenc}
\usepackage{lmodern}
\usepackage{graphicx}
\usepackage{amsmath,amssymb}
\usepackage[a4paper,margin=2cm]{geometry}
\usepackage[absolute,overlay]{textpos}
\setlength{\TPHorizModule}{1mm}
\setlength{\TPVertModule}{1mm}
\usepackage[indonesian]{babel}
\usepackage{hyperref}
\setlength{\emergencystretch}{2em}
% unit: Halaman 32

\begin{document}

% === Halaman 32 (python/native) ===

Efek longsoran (avalanche effect)

• Efek longsoran adalah properti kunci yang diinginkan dari algoritma 

enkripsi,

• Tujuan: menunjukkan perubahan satu bit pada plainteks atau kunci 

menghasilkan perubahan besar pada cipherteks

• Efek longsoran membuat upaya untuk mencoba menebak kunci menjadi

tidak mungkin 

• DES menunjukkan efek longsoran yang kuat

32

\end{document}
